\chapter{2022年全国乙卷(文)}
\section{单选题}

\begin{problem}
已知集合 \(M=\{2,4,6,8,10\}\),\(N=\{x\mid -1<x<6\}\),则 \(M\cap N=\)(\quad)

\choices
    {\(\{2,4\}\)}
    {\(\{2,4,6\}\)}
    {\(\{2,4,6,8\}\)}
    {\(\{2,4,6,8,10\}\)}
\end{problem}

\begin{problem}
设 \((1+2\i)a+b=2\i\),其中 \(a\),\(b\in\R\),则(\quad)

\choices
    {\(a=1\),\(b=-1\)}
    {\(a=1\),\(b=1\)}
    {\(a=-1\),\(b=1\)}
    {\(a=-1\),\(b=-1\)}
\end{problem}

\begin{problem}
已知向量 \(\bs{a}=(2,1)\),\(\bs{b}=(-2,4)\),则 \(\left|\bs{a}-\bs{b}\right|=\)(\quad)

\choices
    {2}
    {3}
    {4}
    {5}
\end{problem}

\begin{problem}
分别统计了甲,乙两位同学 \(16\) 周的各周课外体育运动时长(单位:h),得如下茎叶图:

\begin{center}
\begin{tabular}{r|c|l}
甲 & 茎 & 乙 \\
\hline
\(6\quad1\) & \(5\). & \\
\(8\),\(5\),\(3\),\(0\) & \(6\). & \(3\) \\
\(7\),\(5\),\(3\),\(2\) & \(7\). & \(4\quad6\) \\
\(6\),\(4\),\(2\),\(1\) & \(8\). & \(1\),\(2\),\(2\),\(5\),\(6\),\(6\),\(6\),\(6\) \\
\(4\quad2\) & \(9\). & \(0\),\(2\),\(3\),\(8\) \\
 & \(10\). & \(1\)
\end{tabular}
\end{center}

则下列结论中错误的是(\quad)

\choices
    {甲同学周课外体育运动时长的样本中位数为 \(7.4\)}
    {乙同学周课外体育运动时长的样本平均数大于 \(8\)}
    {甲同学周课外体育运动时长大于 \(8\) 的概率的估计值大于 \(0.4\)}
    {乙同学周课外体育运动时长大于 \(8\) 的概率的估计值大于 \(0.6\)}
\end{problem}

\begin{problem}
若 \(x\),\(y\) 满足约束条件 \(\begin{cases}
x+y\ge2\\
x+2y\le4\\
y\ge0
\end{cases}\),则 \(z=2x-y\) 的最大值是(\quad)

\choices
    {\(-2\)}
    {4}
    {8}
    {12}
\end{problem}

\begin{problem}
设抛物线 \(C:y^2=4x\) 的焦点为 \(F\),点 \(A\) 在 \(C\) 上,点 \(B(3,0)\). 若 \(AF=BF\),则 \(AB=\)(\quad)

\choices
    {2}
    {\(2\sqrt2\)}
    {3}
    {\(3\sqrt2\)}
\end{problem}

\begin{problem}
执行下边的程序框图,输出的 \(n=\)(\quad)

\bitmapfigure[max width=6.0cm]{img/2022/national_paper_b_liberal/q7_fig1.png}

\choices
    {3}
    {4}
    {5}
    {6}
\end{problem}

\begin{problem}
如图是下列四个函数中的某个函数在区间 \([-3,3]\) 的大致图像,则该函数是(\quad)

\choices
    {\(y=\dfrac{-x^3+3x}{x^2+1}\)}
    {\(y=\dfrac{x^3-x}{x^2+1}\)}
    {\(y=\dfrac{2x\cos x}{x^2+1}\)}
    {\(y=\dfrac{2\sin x}{x^2+1}\)}
\end{problem}

\begin{problem}
在正方体 \(ABCD-A_1B_1C_1D_1\) 中,\(E\),\(F\) 分别为 \(AB\),\(BC\) 的中点,则(\quad)

\choices
    {平面 \(B_1EF\perp\) 平面 \(BDD_1\)}
    {平面 \(B_1EF\perp\) 平面 \(A_1BD\)}
    {平面 \(B_1EF//\) 平面 \(A_1AC\)}
    {平面 \(B_1EF//\) 平面 \(A_1C_1D\)}
\end{problem}

\begin{problem}
已知等比数列 \(\{a_n\}\) 的前 \(3\) 项和为 \(168\),\(a_2-a_5=42\),则 \(a_6=\)(\quad)

\choices
    {14}
    {12}
    {6}
    {3}
\end{problem}

\begin{problem}
函数 \(f(x)=\cos x+(x+1)\sin x+1\) 在区间 \([0,2\pi]\) 上的最小值,最大值分别为(\quad)

\choices
    {\(-\dfrac{\pi}{2}\),\(\dfrac{\pi}{2}\)}
    {\(-\dfrac{3\pi}{2}\),\(\dfrac{\pi}{2}\)}
    {\(-\dfrac{\pi}{2}\),\(\dfrac{\pi}{2}+2\)}
    {\(-\dfrac{3\pi}{2}\),\(\dfrac{\pi}{2}+2\)}
\end{problem}

\begin{problem}
已知球 \(O\) 的半径为 \(1\),四棱锥的顶点为 \(O\),底面的四个顶点均在球 \(O\) 的球面上,则当该四棱锥的体积最大时,其高为(\quad)

\choices
    {\(\dfrac{1}{3}\)}
    {\(\dfrac{1}{2}\)}
    {\(\dfrac{\sqrt3}{3}\)}
    {\(\dfrac{\sqrt2}{2}\)}
\end{problem}

\section{填空题}

\begin{problem}
记 \(S_n\) 为等差数列 \(\{a_n\}\) 的前 \(n\) 项和. 若 \(2S_3=3S_2+6\),则公差 \(d=\fillinblank{}\).
\end{problem}

\begin{problem}
从甲,乙等 \(5\) 名同学中随机选 \(3\) 名参加社区服务工作,则甲,乙都入选的概率为 \(\fillinblank{}\).
\end{problem}

\begin{problem}
过四点 \((0,0)\),\((4,0)\),\((-1,1)\),\((4,2)\) 中任意三点可得的一个圆的方程为 \(\fillinblank{}\).
\end{problem}

\begin{problem}
若 \(f(x)=\ln\left|a+\dfrac{1}{1-x}\right|+b\) 是奇函数,则 \(a=\fillinblank{}\),\(b=\fillinblank{}\).
\end{problem}

\section{解答题}

\begin{problem}
已知 \(\triangle ABC\) 的内角 \(A\),\(B\),\(C\) 的对边分别为 \(a\),\(b\),\(c\),且 \(\sin C\sin\left(A-B\right)=\sin B\sin\left(C-A\right)\).

\begin{enumerate}
\item 若 \(A=2B\),求 \(C\);
\item 证明:\(2a^2=b^2+c^2\).
\end{enumerate}
\end{problem}

\begin{problem}
如图,四面体 \(ABCD\) 中,\(AD\perp CD\),\(AD=CD\),\(\angle ADB=\angle CDB\),\(E\) 为 \(AC\) 的中点.

\begin{center}
\bitmapfigure[max width=6.8cm]{img/2022/national_paper_B_liberal/q18_fig1.png}
\end{center}

\circled{1} 证明:平面 \(BED\perp\) 平面 \(ACD\);

\circled{2} 设 \(AB=BD=2\),\(\angle ACB=60^\circ\),点 \(F\) 在 \(BD\) 上,当 \(\triangle AFC\) 的面积最小时,求三棱锥 \(F-ABC\) 的体积.
\end{problem}

\begin{problem}
某地经过多年的环境治理,已将荒山改造成了绿水青山. 为估计一林区某种树木的总材积量,随机选取了 \(10\) 棵这种树木,测量每棵树的根部横截面积(单位:\(\text{m}^2\))和材积量(单位:\(\text{m}^3\)),得到如下数据:

\begin{center}
\begin{tabular}{c|ccccc}
样本号 \(i\) & 1 & 2 & 3 & 4 & 5 \\\hline
根部横截面积 \(x_i\) & 0.04 & 0.06 & 0.04 & 0.08 & 0.08 \\
材积量 \(y_i\) & 0.25 & 0.40 & 0.22 & 0.54 & 0.51 \\
\end{tabular}

\vspace{0.5em}

\begin{tabular}{c|cccc|c}
样本号 \(i\) & 6 & 7 & 8 & 9 & 10 \\\hline
根部横截面积 \(x_i\) & 0.05 & 0.05 & 0.07 & 0.07 & 0.06 \\
材积量 \(y_i\) & 0.34 & 0.36 & 0.46 & 0.42 & 0.40 \\
\end{tabular}
\end{center}

并计算得 \(\sum_{i=1}^{10}x_i=0.6\),\(\sum_{i=1}^{10}y_i=3.9\),\(\sum_{i=1}^{10}x_i^2=0.038\);
\(\sum_{i=1}^{10}y_i^2=1.6158\),\(\sum_{i=1}^{10}x_iy_i=0.2474\).

\begin{enumerate}
\item 估计该林区这种树木平均一棵的根部横截面积与平均一棵的材积量;
\item 求该林区这种树木的根部横截面积与材积量的样本相关系数(精确到 \(0.01\));
\item 现测量了该林区所有这种树木的根部横截面积,并得到所有这种树木的根部横截面积总和为 \(186\text{ m}^2\). 已知树木的材积量与其根部横截面积近似成正比,利用以上数据给出该林区这种树木的总材积量的估计值.
\end{enumerate}

附:相关系数 \(r=\dfrac{\sum_{i=1}^n(x_i-\bar{x})(y_i-\bar{y})}{\sqrt{\sum_{i=1}^n(x_i-\bar{x})^2\sum_{i=1}^n(y_i-\bar{y})^2}}\),且 \(\sqrt{1.896}\approx1.377\).
\end{problem}

\begin{problem}
已知函数 \(f(x)=ax-\dfrac{1}{x}-(a+1)\ln x\),\(x>0\).

\begin{enumerate}
    \item 当 \(a=0\) 时,求 \(f(x)\) 的最大值;
    \item 若 \(f(x)\) 恰有一个零点,求 \(a\) 的取值范围.
\end{enumerate}
\end{problem}

\begin{problem}
已知椭圆 \(E\) 的中心为坐标原点,对称轴为 \(x\) 轴,\(y\) 轴,且过 \(A(0,-2)\),\(B\left(\dfrac32,-1\right)\) 两点.

\begin{enumerate}
\item 求 \(E\) 的方程;

\item 设过点 \(P(1,-2)\) 的直线交 \(E\) 于 \(M\),\(N\) 两点,过 \(M\) 且平行于 \(x\) 轴的直线与线段 \(AB\) 交于点 \(T\),点 \(H\) 满足 \(MT=TH\). 证明:直线 \(HN\) 过定点.
\end{enumerate}
\end{problem}

\section{选考题:解答题}

\begin{problem}
在直角坐标系 \(xOy\) 中,曲线 \(C\) 的参数方程为 \(x=\sqrt3\cos2t\),\(y=2\sin t\)(\(t\) 为参数),以坐标原点为极点,\(x\) 轴正半轴为极轴建立极坐标系,已知直线 \(l\) 的极坐标方程为 \(\rho\sin\left(\theta+\frac{\pi}{3}\right)+m=0\).

\begin{enumerate}
\item 写出 \(l\) 的直角坐标方程;
\item 若 \(l\) 与 C 有公共点,求 \(m\) 的取值范围.
\end{enumerate}
\end{problem}

\begin{problem}
已知 \(a\),\(b\),\(c>0\),且 \(a^{\frac32}+b^{\frac32}+c^{\frac32}=1\). 证明:

\begin{enumerate}
\item \(abc\le\dfrac19\);
\item \(\dfrac{a}{b+c}+\dfrac{b}{c+a}+\dfrac{c}{a+b}\le\dfrac{1}{2\sqrt{abc}}\).
\end{enumerate}
\end{problem}

